\chapter{The Three-Equation Model}\label{ch:b2:three-equation-model}

Central banks set an interest rate, and they set it by looking at inflation and at how far output is from its normal
level. This chapter builds the model that describes that practice in its simplest form: a supply side where wage and price
setting fix equilibrium unemployment, a demand curve, a Phillips curve and a central bank that chooses the interest rate.
It is a toy model, small enough to solve by hand. The chart compares the American policy rate with a simple rule published
in 1993.

\begin{openingchart}{United States: the federal funds rate and the rule $r = p + 0.5y + 0.5(p - 2) + 2$ of Taylor (1993), quarterly, 1960--2008; $p$ is CPI inflation over four quarters,
  $y$ the per cent deviation of real GDP from a Hodrick--Prescott trend.}{Federal Reserve, H.15 (RIFSPFF\_N.M); statsmodels
  macrodata (cpi, realgdp); our computation}
\begin{tikzpicture}
\begin{axis}[omrecord, width=0.72\linewidth, height=0.42\linewidth, xmin=1960, xmax=2009, ymin=-2, ymax=25,
    xticklabel style={/pgf/number format/1000 sep=}, ylabel={Per cent},
    legend style={at={(0.97,0.97)}, anchor=north east, font=\scriptsize}]
  \addplot[omDef, thick] table[x=t, y=ffr] {figdata/bachelor-2/27-three-equation-model/out/rule.csv};
  \addlegendentry{federal funds rate}
  \addplot[omThm, thick, dashed] table[x=t, y=rule] {figdata/bachelor-2/27-three-equation-model/out/rule.csv};
  \addlegendentry{1993 rule}
\end{axis}
\end{tikzpicture}
\end{openingchart}

% figdata: rule_dec (mean deviation: 1960s -0.49, 1970s -4.74, 1980s 0.88, 1990s -0.21, 2000s -2.13); rule_summary (MAD 1960-79 2.90, 1987-2007 1.40)
In the 1970s the funds rate was on average 4.74 points below the rule; from 1987 to 2007 it stayed on average 1.40 points from it,
against 2.90 in 1960--79. The rule was written in 1993 to describe policy since 1987, so the fit after 1987 is partly by
construction; the 1970s gap is not. After 2001, and again after 2007, the rate went well below the rule.

\begin{recall}
From the Year 1 volume: the IS curve and an interest-rate rule for monetary policy; the natural rate of unemployment;
aggregate supply and the expectations-augmented Phillips curve; the sacrifice ratio. From
\cref{ch:b2:expectations-and-lucas-critique}: the \omterm{def:b2:expectations-and-lucas-critique:loss}{central-bank loss function}, \omterm{def:b2:expectations-and-lucas-critique:rational-expectations}{rational expectations} and the \omterm{def:b2:expectations-and-lucas-critique:lucas-critique}{Lucas critique}.
\end{recall}

\begin{notation}[Symbols of this chapter]\label{not:b2:three-equation-model:symbols}
$x$ is the output gap, the per cent deviation of output from its equilibrium level (a local symbol here); $\pi$ is inflation and $\pi^{\mathrm T}$ the target; $r$ is the
real interest rate, $i$ the nominal rate, $r^*$ the real rate that keeps output at equilibrium; $a$ is the effect of the real rate on
output, $\alpha$ the slope of the Phillips curve, $\lambda$ the bank's weight on output; $\mu$ is the price markup as a share of output per worker.
\end{notation}

\section{Wage setting and price setting}\label{sec:b2:three-equation-model:ws-ps}

\begin{definition}[Wage-setting curve]\label{def:b2:three-equation-model:ws}
The \emph{wage-setting curve}\index{wage-setting curve} gives the real wage that results from wage setting, by bargaining or by
employers paying efficiency wages, at each unemployment rate; it falls as unemployment rises, since workers' outside
options worsen.
\end{definition}

\begin{definition}[Price-setting curve]\label{def:b2:three-equation-model:ps}
The \emph{price-setting curve}\index{price-setting curve} gives the real wage implied by firms' pricing: with prices set as a
markup on unit labour cost, the real wage equals output per worker times $1 - \mu$, independent of unemployment when output per worker is
constant.
\end{definition}

Where the two curves cross, the real wage workers obtain is the real wage firms' prices leave them. At any lower
unemployment, wage setters push for more than price setters allow: wages chase prices and prices chase wages, and inflation
rises. At higher unemployment it falls. The crossing is equilibrium unemployment, and output there is equilibrium output.

\begin{omfigure}
\begin{tikzpicture}[econ]
\begin{axis}[omrecord, grid=none, xmin=0, xmax=0.12, ymin=0.6, ymax=1.15, width=0.55\linewidth, height=0.4\linewidth,
    xlabel={Unemployment rate, $u$}, ylabel={Real wage}, scaled x ticks=false, xticklabel style={/pgf/number format/fixed},
    legend style={at={(1.03,0.5)}, anchor=west, font=\small}]
  \addplot[omDef, very thick, domain=0:0.1, samples=2] {1.1 - 5*x};
  \addlegendentry{WS: $1.1 - 5u$}
  \addplot[omThm, very thick, domain=0:0.12, samples=2] {0.8};
  \addlegendentry{PS, $\mu = 0.20$}
  \addplot[omThm, thick, dashed, domain=0:0.12, samples=2] {0.75};
  \addlegendentry{PS, $\mu = 0.25$}
  \node[eqpoint] at (axis cs:0.06,0.8) {};
  \node[eqpoint] at (axis cs:0.07,0.75) {};
  \draw[guide] (axis cs:0.06,0.8) -- (axis cs:0.06,0.6);
  \draw[guide] (axis cs:0.07,0.75) -- (axis cs:0.07,0.6);
\end{axis}
\end{tikzpicture}

{\small Brenland's wage-setting and \omterm{def:b2:three-equation-model:ps}{price-setting curves} (the weekend problem). Equilibrium unemployment is 6 per cent; a rise
in the markup from 0.20 to 0.25 lowers the price-setting real wage and raises it to 7 per cent.}
\end{omfigure}

\section{A demand curve with expectations}\label{sec:b2:three-equation-model:is}

Spending depends on the real interest rate, and with a lag: investment plans and durable purchases take time to respond.
In gap form, $x_t = -a(r_{t-1} - r^*)$. The rate $r^*$ is the real rate at which output is at its equilibrium level; it changes when the demand curve shifts. The
central bank sets a nominal rate, and the real rate is the nominal rate minus the inflation people expect.

\section{A Phillips curve}\label{sec:b2:three-equation-model:pc}

With expectations of inflation equal to last period's inflation, the Phillips curve is $\pi_t = \pi_{t-1} + \alpha x_t$: inflation rises when output is above its
equilibrium level, the gap between the curves of the WS--PS diagram. Each Phillips curve is drawn for one value of past
inflation, and it shifts each period.

\section{A monetary rule}\label{sec:b2:three-equation-model:mr}

\begin{definition}[Three-equation model]\label{def:b2:three-equation-model:three-equation}
The \emph{three-equation model}\index{three-equation model} of monetary policy is the model made of a demand (IS)
equation, a Phillips curve and a monetary rule derived from the central bank's loss function.
\end{definition}

The bank's loss is $L = (\pi_t - \pi^{\mathrm T})^2 + \lambda x_t^2$. Because of the lag in the IS curve, it cannot affect this period's output: it sets $r_{t-1}$ now to choose next period's point on next
period's Phillips curve. Minimising the loss subject to $\pi_t = \pi_{t-1} + \alpha x_t$: $2\alpha(\pi_t - \pi^{\mathrm T}) + 2\lambda x_t = 0$.

\begin{definition}[Monetary-rule curve]\label{def:b2:three-equation-model:mr}
The \emph{monetary-rule curve}\index{monetary-rule curve} is the set of combinations of inflation and output gap that the
central bank chooses, given its loss function and the Phillips curve it faces: here $x = -\dfrac{\alpha}{\lambda}(\pi - \pi^{\mathrm T})$.
\end{definition}

Combining the MR curve with the Phillips curve and the IS curve gives the interest rate the bank must set. With $\pi_{t-1}$ known,
\begin{equation}\label{eq:b2:three-equation-model:rule}
\pi_t - \pi^{\mathrm T} = \frac{\pi_{t-1} - \pi^{\mathrm T}}{1 + \alpha^2/\lambda}, \qquad r_{t-1} - r^* = \frac{\alpha}{a(\lambda + \alpha^2)}(\pi_{t-1} - \pi^{\mathrm T}).
\end{equation}
The bank raises the real rate in proportion to the inflation gap, and each period inflation closes the same fraction of its
gap. The output gap does not appear in the rule only because the model has no demand shocks that the bank can see
coming; with them, it appears too.

\begin{definition}[Taylor rule]\label{def:b2:three-equation-model:taylor-rule}
A \emph{Taylor rule}\index{Taylor rule} is an interest-rate rule in which the central bank sets its policy rate as a linear
function of inflation, or its gap from the target, and of the output gap, $i = r^* + \pi + \phi_\pi(\pi - \pi^{\mathrm T}) + \phi_x x$.
\end{definition}

\begin{record}[Rules and rates, 1960--2007]
% ledger: b2:taylor1993, b2:cgg2000; figdata: rule_summary
Taylor (1993) proposed the rule of the opening chart, with an inflation target of 2 per cent and an equilibrium real rate of 2 per cent,
and found that it fitted actual policy remarkably well over 1987--1992, except in 1987, when the Federal Reserve eased after the
stock-market crash. Clarida, Gal\'i and Gertler (2000) estimated forward-looking rules for the post-war United States and found that
policy after 1979 was much more sensitive to expected inflation than before, and that the later rule was stabilising in a
simple model, the earlier one not. In the chart the rate was on average 4.74 points below the rule in the 1970s.
\end{record}

\section{Shocks and responses}\label{sec:b2:three-equation-model:shocks}

Take Brenland's parameters: $\alpha = 0.5$, $\lambda = 0.5$, $a = 0.5$, $r^* = 2$ and a target of 2 per cent. An inflation shock raises inflation to 5 per cent. The MR curve is $x = -(\pi - 2)$, and the bank aims at the point
of next period's Phillips curve $\pi = 5 + 0.5x$ that lies on it: $x = -2$, $\pi = 4$. To get there it sets the real rate at $r = 2 + 2/0.5 = 6$ per cent. Next period, with inflation 4, the
Phillips curve has shifted down, and so on.

\begin{omfigure}
\begin{tikzpicture}[econ]
\begin{axis}[omrecord, grid=none, xmin=-3, xmax=1.5, ymin=1, ymax=6, width=0.55\linewidth, height=0.45\linewidth,
    xlabel={Output gap, $x$}, ylabel={Inflation, $\pi$}, clip=false]
  \draw[black!40] (axis cs:0,1) -- (axis cs:0,6);
  \addplot[omThm, very thick, domain=-3:1, samples=2] {2 - x};
  \addplot[omDef, thick, domain=-3:1.5, samples=2] {5 + 0.5*x};
  \addplot[omDef, thick, dashed, domain=-3:1.5, samples=2] {4 + 0.5*x};
  \addplot[omDef, thick, dotted, domain=-3:1.5, samples=2] {3.3333 + 0.5*x};
  \node[eqpoint, label={[font=\small]right:A}] at (axis cs:0,5) {};
  \node[eqpoint, label={[font=\small]below left:B}] at (axis cs:-2,4) {};
  \node[eqpoint, label={[font=\small]below left:C}] at (axis cs:-1.3333,3.3333) {};
  \node[eqpoint, label={[font=\small]right:Z}] at (axis cs:0,2) {};
  \node[font=\small, anchor=south west] at (axis cs:-2.95,5.0) {MR};
  \node[font=\scriptsize, anchor=west] at (axis cs:1.55,5.75) {PC, $\pi_{-1} = 5$};
  \node[font=\scriptsize, anchor=west] at (axis cs:1.55,4.75) {PC, $\pi_{-1} = 4$};
  \node[font=\scriptsize, anchor=west] at (axis cs:1.55,4.08) {PC, $\pi_{-1} = 3.33$};
\end{axis}
\end{tikzpicture}

{\small Adjustment after an inflation shock. The economy starts at A, with inflation 5 and output at equilibrium. The bank
chooses B, where the Phillips curve for past inflation 5 meets the MR curve; then C on the next Phillips curve (dashed); and so on down
the MR curve to Z, the target with a zero gap.}
\end{omfigure}

The listing follows the path period by period. The nominal rate is the real rate plus expected inflation, which here is
current inflation.

\omcode{code/bachelor-2/27-three-equation-model/listings/three_eq.py}{1}{27}{The path of inflation and interest rates
after an inflation shock.}\label{lst:b2:three-equation-model:path}

\omoutput{code/bachelor-2/27-three-equation-model/outputs/three_eq.txt}{Inflation closes a third of its gap each period;
the real rate starts 4 points above $r^*$; the cumulated output loss is twice the fall in inflation.}\label{out:b2:three-equation-model:path}

\begin{omfigure}
\begin{tikzpicture}[econ]
\begin{axis}[omrecord, grid=none, xmin=-0.3, xmax=8.3, ymin=0, ymax=12, width=0.55\linewidth, height=0.38\linewidth,
    xlabel={Period}, ylabel={Per cent}, legend style={at={(1.03,0.5)}, anchor=west, font=\small}]
  \addplot[omDef, thick, mark=*, mark size=1.4pt, samples at={0,...,8}] {2 + 4*(2/3)^x + 2 + 3*(2/3)^x};
  \addlegendentry{nominal rate}
  \addplot[omThm, thick, mark=square*, mark size=1.4pt, samples at={0,...,8}] {2 + 4*(2/3)^x};
  \addlegendentry{real rate}
  \addplot[omMeth, thick, mark=triangle*, mark size=1.6pt, samples at={0,...,8}] {2 + 3*(2/3)^x};
  \addlegendentry{inflation}
\end{axis}
\end{tikzpicture}

{\small The interest-rate path of the listing: both rates jump and then fall back with inflation. The nominal rate moves by
more than inflation: it starts 7 points above its long-run level of 4 when inflation is 3 points above target.}
\end{omfigure}

% code: listings/three_eq.py; exercises/solutions.py cumulated_loss
The cumulated output loss is 6 per cent of a period's output, twice the 3-point fall in inflation: a sacrifice ratio of 2, which is $1/\alpha$. It does not depend on $\lambda$: a
bank that cares more about output spreads the same total loss over more periods. With a Phillips curve built on past
inflation, each point of disinflation must be bought with $1/\alpha$ points of output below equilibrium.

\begin{model}{The three-equation model}
\assumes A Phillips curve with expectations equal to past inflation; demand that responds to the real rate with a lag; a
central bank that minimises a quadratic loss and knows the economy's parameters; no financial sector.
\predicts An interest-rate rule that raises the real rate with the inflation gap; inflation and output returning to
target along the MR curve; a sacrifice ratio of $1/\alpha$.
\tested From the mid-1980s to 2007 the American funds rate stayed close to a simple rule of this form (Taylor 1993; the
opening chart); policy after 1979 responded more strongly to inflation than before (Clarida, Gal\'i and Gertler 2000).
\breaks In 2008 a financial crisis, outside the model, drove rates to zero; estimates on American states find a Phillips
curve that is flat (Hazell, Herre\~no, Nakamura and Steinsson 2022), so the output cost of moving inflation would be large; in 2021--22 inflation rose a
long way before rates did.
\end{model}

\begin{omchart}{United States: CPI inflation over twelve months and the federal funds rate, monthly, January 2019 to
  October 2026, per cent.}{Bureau of Labor Statistics, CPI-U; Federal Reserve, H.15 (RIFSPFF\_N.M); our computation}
\begin{tikzpicture}
\begin{axis}[omrecord, width=0.7\linewidth, height=0.38\linewidth, xmin=2019, xmax=2026.9, ymin=0, ymax=10,
    xticklabel style={/pgf/number format/1000 sep=}, ylabel={Per cent},
    legend style={at={(0.03,0.97)}, anchor=north west, font=\scriptsize}]
  \addplot[omMeth, thick] table[x=t, y=infl] {figdata/bachelor-2/27-three-equation-model/out/recent.csv};
  \addlegendentry{CPI inflation}
  \addplot[omDef, thick] table[x=t, y=ffr] {figdata/bachelor-2/27-three-equation-model/out/recent.csv};
  \addlegendentry{federal funds rate}
\end{axis}
\end{tikzpicture}
\end{omchart}

\begin{record}[2021--23]
% figdata: recent_summary (CPI peak 9.1 in 2022-06 with ffr 1.21; ffr first above 0.5 in 2022-05; ffr peak 5.33 in 2023-08; 2026-10: CPI 3.9, ffr 3.88)
American CPI inflation over twelve months peaked at 9.1 per cent in June 2022. The funds rate first rose above half a point in
May 2022 and was 1.21 per cent in the month of the peak; it reached 5.33 per cent in August 2023. In the model the bank responds
to inflation at once and the real rate rises with it; in 2021--22 the real rate, measured with current inflation, fell to
far below zero before it rose. Whether this was a misreading of the shocks as temporary or a deliberate choice is
outside the model; the record is that inflation rose a long way before policy moved.
\end{record}

\section{The rule must react more than one for one}\label{sec:b2:three-equation-model:taylor-principle}

\begin{definition}[Taylor principle]\label{def:b2:three-equation-model:taylor-principle}
The \emph{Taylor principle}\index{Taylor principle} is the requirement that the nominal policy rate rise by more than one
point for each point of rise in inflation, so that the real rate rises when inflation rises.
\end{definition}

\begin{proposition}[The Taylor principle in the three-equation model]\label{prop:b2:three-equation-model:principle}
Let the bank set $i_t = r^* + \pi^{\mathrm T} + \phi(\pi_t - \pi^{\mathrm T})$, with expected inflation $\pi_t$. Inflation returns to target from any starting point if and only if $1 < \phi < 1 + 2/(a\alpha)$.
\end{proposition}

\begin{proof}
The real rate is $r_t = i_t - \pi_t = r^* + (\phi - 1)(\pi_t - \pi^{\mathrm T})$. Then $x_{t+1} = -a(\phi - 1)(\pi_t - \pi^{\mathrm T})$, and the Phillips curve gives $\pi_{t+1} - \pi^{\mathrm T} = (1 - a\alpha(\phi - 1))(\pi_t - \pi^{\mathrm T})$. The gap shrinks to zero if and only if $|1 - a\alpha(\phi - 1)| < 1$, that is
$0 < a\alpha(\phi - 1) < 2$.
\end{proof}

With $\phi < 1$, a rise in inflation lowers the real rate, raises output and raises inflation further. With Brenland's parameters the derived rule
has $\phi = 1 + 4/3$: the nominal rate rises by 2.33 points per point of inflation, and each period the gap shrinks to $1 - 0.25 \times 4/3 = 2/3$ of its value. A rule with
$\phi = 0.5$ would multiply the gap by 1.125 each period.

\section{The zero lower bound}\label{sec:b2:three-equation-model:zlb}

Cash pays a nominal return of zero, so the nominal rate cannot go much below zero. When inflation is low and $r^*$ has fallen, the rule
may call for a negative nominal rate. With $r^* = 1$, inflation 0.5 and Brenland's rule, the desired real rate is $1 - \tfrac43 \times 1.5 = -1$, a nominal rate of $-0.5$. Set at zero, the real rate is only $-0.5$: next
period's gap is 0.75 instead of the 1.0 the bank wanted, and inflation returns more slowly.

\begin{omfigure}
\begin{tikzpicture}[econ]
\begin{axis}[omrecord, grid=none, xmin=-1, xmax=4, ymin=-3, ymax=8, width=0.5\linewidth, height=0.4\linewidth,
    xlabel={Inflation, $\pi$}, ylabel={Nominal rate, $i$}]
  \draw[black!40] (axis cs:-1,0) -- (axis cs:4,0);
  \addplot[omDef, thick, dashed, domain=-1:0.714, samples=2] {7/3*x - 5/3};
  \addplot[omDef, very thick, domain=0.714:4, samples=2] {7/3*x - 5/3};
  \addplot[omDef, very thick, domain=-1:0.714, samples=2] {0};
  \node[eqpoint] at (axis cs:0.714,0) {};
\end{axis}
\end{tikzpicture}

{\small The rule $i = r^* + \pi + \tfrac43(\pi - 2)$ with $r^* = 1$, and the zero lower bound: below an inflation rate of about 0.71 per cent, the rule asks for a negative rate
(dashed), and the rate set stays at zero.}
\end{omfigure}

The model's forward-looking variant replaces past inflation by expected future inflation in the Phillips curve. Then a
promise to keep rates low for longer, after the bound stops binding, raises expected inflation and lowers the real rate
today; the promise works only if it is believed, which is the credibility problem of
\cref{ch:b2:expectations-and-lucas-critique}.

\section{Exercises}\label{sec:b2:three-equation-model:exercises}

\begin{exercise}[$\star$]\label{exo:b2:three-equation-model:1}
Read the opening chart. In which decade was the funds rate furthest below the rule on average? Compare the average distance
from the rule in 1960--79 and 1987--2007.
\end{exercise}

\begin{exercise}[$\star$]\label{exo:b2:three-equation-model:2}
Read the chart of 2019--2026. When did inflation peak, and what was the funds rate then? When did the rate first exceed half a
point, and when did it peak?
\end{exercise}

\begin{exercise}[$\star$]\label{exo:b2:three-equation-model:3}
The \omterm{def:b2:three-equation-model:ws}{wage-setting curve} is $1.2 - 4u$ and the price-setting real wage 0.8. Find equilibrium unemployment.
\end{exercise}

\begin{exercise}[$\star$]\label{exo:b2:three-equation-model:4}
With inflation at 4 per cent and output 1 per cent below trend, what funds rate does the 1993 rule give?
\end{exercise}

\begin{exercise}[$\star\star$]\label{exo:b2:three-equation-model:5}
With $\alpha = \lambda = a = 1$, find the MR curve, the fraction of the inflation gap left after one period, and the coefficient of the interest-rate rule.
\end{exercise}

\begin{exercise}[$\star\star$]\label{exo:b2:three-equation-model:6}
Predict, then compute: with Brenland's parameters, a one-period demand shock raises the output gap to 2 with past inflation at target.
What are inflation this period, inflation and the gap the bank chooses for next period, and the real and nominal rates it sets?
\end{exercise}

\begin{exercise}[$\star\star$]\label{exo:b2:three-equation-model:7}
With $r^* = 1$, inflation 0.5 and Brenland's rule, find the nominal rate the rule asks for, the rate set, and next period's output gap with
and without the bound.
\end{exercise}

\begin{exercise}[$\star\star$]\label{exo:b2:three-equation-model:8}
Modify the listing to set $\lambda = 2$. What happens to the first period's inflation, real rate and output gap, and to the cumulated output loss?
\end{exercise}

\begin{exercise}[$\star\star$]\label{exo:b2:three-equation-model:9}
\textbf{(D)} With the macrodata and H.15 series, compute the funds rate minus the 1993 rule, quarterly 1960--2008, and report
its mean and its mean absolute value by decade.
\end{exercise}

\begin{exercise}[$\star\star\star$]\label{exo:b2:three-equation-model:10}
Derive \eqref{eq:b2:three-equation-model:rule} from the loss function, the Phillips curve and the IS curve.
\end{exercise}

\begin{exercise}[$\star\star\star$]\label{exo:b2:three-equation-model:11}
Show that the derived rule satisfies the \omterm{def:b2:three-equation-model:taylor-principle}{Taylor principle} for all positive $a$, $\alpha$ and $\lambda$, and find the $\phi$ that the proposition's upper bound would allow with
Brenland's parameters.
\end{exercise}

\begin{exercise}[$\star\star\star$]\label{exo:b2:three-equation-model:12}
Show that, after an inflation shock that leaves a gap $g_0$, the cumulated output loss along the MR path is $g_0/\alpha$ whatever $\lambda$.
\end{exercise}

\section{Problem: Brenland's central bank meets a supply shock}\label{sec:b2:three-equation-model:problem}

\begin{problem}[{Weekend problem --- Brenland's central bank meets a supply shock: wage and price setting, the Phillips and
MR curves, the path of rates, and the cost of disinflation}]
\label{pb:b2:three-equation-model:1}
Brenland's \omterm{def:b2:three-equation-model:ws}{wage-setting curve} is $1.1 - 5u$ and firms price at a markup of 0.20 of output per worker, which is 1. Demand is $x_t = -0.5(r_{t-1} - 2)$; the Phillips curve is
$\pi_t = \pi_{t-1} + 0.5x_t$; the central bank's loss is $(\pi - 2)^2 + 0.5x^2$. A supply shock has just raised inflation to 5 per cent, with output at equilibrium.

\medskip
\noindent\textbf{Part I --- Wage and price setting.}
\begin{enumerate}
  \item What does the \omterm{def:b2:three-equation-model:ws}{wage-setting curve} say about bargaining?
  \item Find the price-setting real wage.
  \item Find equilibrium unemployment.
  \item What happens to inflation if unemployment is held at 4 per cent?
  \item The markup rises to 0.25. Find the new equilibrium unemployment.
\end{enumerate}

\medskip
\noindent\textbf{Part II --- The Phillips and MR curves.}
\begin{enumerate}[resume]
  \item Write the Phillips curve the bank faces for next period.
  \item Derive the MR curve.
  \item Find the point the bank chooses for next period.
  \item What real rate does it set now?
  \item What nominal rate, if expected inflation is current inflation?
  \item Write the interest-rate rule as $r - r^*$ in terms of $\pi - \pi^{\mathrm T}$.
\end{enumerate}

\medskip
\noindent\textbf{Part III --- The path.}
\begin{enumerate}[resume]
  \item What fraction of the inflation gap is left each period?
  \item Find inflation and the output gap in period 2.
  \item Find the real rate set in period 1.
  \item After how many periods is inflation below 2.5 per cent?
  \item Check the first rows of the listing's output against your answers.
  \item What nominal-rate coefficient on inflation does the rule imply? Does it satisfy the \omterm{def:b2:three-equation-model:taylor-principle}{Taylor principle}?
\end{enumerate}

\medskip
\noindent\textbf{Part IV --- The cost.}
\begin{enumerate}[resume]
  \item Find the cumulated output loss.
  \item Find the sacrifice ratio.
  \item With a weight of 2 on output, find the first period's inflation and output gap.
  \item What is the cumulated loss with that weight?
  \item Why is the total cost the same, and what does the weight change?
  \item What would a flatter Phillips curve, $\alpha = 0.25$, do to the sacrifice ratio?
  \item If the shock had come with a fall in $r^*$ to 0.5, would the zero lower bound have mattered at the start? Why not?
  \item State the result: the interest-rate path, the cumulated output loss and the sacrifice ratio of the disinflation.
\end{enumerate}
\end{problem}
